\begin{resumo}[Abstract]
\begin{otherlanguage*}{english}
\sisetup{output-decimal-marker = {.}}
ISO 20816-3 evaluates the vibration of industrial machines by the RMS velocity measured on the housing and by changes in that value, and a consumer micro-electro-mechanical system (MEMS) accelerometer costs one to two orders of magnitude less than the piezoelectric accelerometer used as the reference for this measurement. In two published low-cost nodes that use the Isolation Forest detector, the false alarm rate was \SI{22}{\percent} to \SI{38}{\percent} of normal operation, and none of the works reviewed assessed the false alarm rate between measurement rounds on the same machine in an industrial plant. This work assessed whether a low-cost acquisition system that transmits the raw vibration signal can detect, with a model trained only on each machine's normal operation, changes of \SI{0,7}{\milli\meter\per\second}, \SI{25}{\percent} of the upper limit of the standard's acceptable vibration zone, with a limited false alarm rate between rounds. A transmitter with an ESP32-C3 and an MPU-6050 module acquired \SI{0,512}{\second} windows at \SI{1}{\kilo\hertz} and sent them over LoRa radio to a gateway, which forwarded them to an application programming interface, and the system was verified by automated tests and end-to-end bench tests. Including taxes, the transmitter and gateway components cost R\$~\num{176,04}, excluding the battery. On 2 November 2026, it acquired \num{\campoJanelasOperacao} windows in \num{\campoEnsaios} tests on \num{\campoMaquinas} machines of an industrial plant, all in normal operation. Four hypotheses were evaluated against criteria registered before any calculation and after exploratory analyses, the fourth by amendment: false alarm between two rounds on the same machine (H1), smallest detectable change in RMS velocity below \SI{0,7}{\milli\meter\per\second} (H2), detection of an unbalance signature digitally added to real windows (H3), and detection of real unbalance in a public laboratory dataset (H4). The analysed windows arrived complete, and no hypothesis was corroborated. H1 was inconclusive, with \num{\hUmCalandraIFDireto}~\% of the second-round windows of the plate rolling machine flagged, and H2 was refuted, with a smallest detectable change of up to \num{\hDoisMmdCalandraDois}~mm/s. H3 was refuted, with at most \num{\hTresFracaoMaxCalandra}~\% of the windows flagged up to \SI{2,8}{\milli\meter\per\second} of signature, on the plate rolling machine, which already had \num{\apUmFracaoSemAssinaturaCalandra}~\% flagged without the signature. H4 was refuted, with a \num{\hQuatroFracaoZero}~\% false alarm rate after the unbalance holder was reassembled, although \num{\hQuatroFracaoTres}~\% to \num{\hQuatroFracaoQuatro}~\% of the unbalanced windows were flagged. In post hoc analyses, H1 would be corroborated with \num{\apTresVereditoUmCorroborada} of the \num{\apTresSementes} random seeds of the detector, the smallest detectable change was \num{\apCincoRazaoMin} to \num{\apCincoRazaoMax} times that of the sensor at rest, and the Isolation Forest score stayed below the threshold on three of the four machines even with the descriptors changed by the signature outside the training range. In these data, detection was limited by the detector's score, by a baseline from a single round, with \num{\campoJanelasGuilhotinaUm} to \num{\campoJanelasCalandraUm} windows per machine, and by the window-to-window variation measured on the machines, which combines the machine's vibration and the response of the sensor mounting.

\vspace{\onelineskip}
\noindent
\textbf{Keywords}: fault detection; predictive maintenance; vibration analysis; MEMS accelerometer; Isolation Forest.
\end{otherlanguage*}
\end{resumo}
